\documentclass[11pt]{article}

\usepackage[T1]{fontenc}
\usepackage[utf8]{inputenc}
\usepackage{lmodern}
\usepackage{microtype}
\usepackage[a4paper,margin=1in]{geometry}
\usepackage{amsmath,amssymb,amsthm,mathtools,mathrsfs}
\usepackage{booktabs,array,longtable}
\usepackage{enumitem}
\usepackage{xcolor}
\usepackage{hyperref}
\usepackage{listings}
\usepackage{fancyhdr}

\definecolor{linkblue}{RGB}{25,75,135}
\definecolor{codegray}{RGB}{247,248,250}
\hypersetup{
  colorlinks=true,
  linkcolor=linkblue,
  citecolor=linkblue,
  urlcolor=linkblue,
  pdftitle={Exact quarter-base extrapolation for dyadic samples of Rvachev's up-function spline prefixes},
  pdfauthor={Technical note accompanying the ProveIt Fabius-function development},
  pdfsubject={Rvachev up-function, finite sinc products, dyadic rational values, q-binomial extrapolation},
  pdfkeywords={Rvachev up-function, Fabius function, dyadic rational, q-binomial, q-Pochhammer, Richardson extrapolation}
}

\pagestyle{fancy}
\fancyhf{}
\fancyhead[C]{\small\nouppercase{\leftmark}}
\fancyfoot[C]{\thepage}
\setlength{\headheight}{14pt}

\setlist[itemize]{leftmargin=2em,itemsep=0.2em,topsep=0.4em}
\setlist[enumerate]{leftmargin=2.4em,itemsep=0.25em,topsep=0.4em}

\lstset{
  basicstyle=\ttfamily\small,
  backgroundcolor=\color{codegray},
  frame=single,
  rulecolor=\color{black!20},
  breaklines=true,
  columns=fullflexible,
  showstringspaces=false,
  keepspaces=true
}

\newtheorem{theorem}{Theorem}[section]
\newtheorem{lemma}[theorem]{Lemma}
\newtheorem{proposition}[theorem]{Proposition}
\newtheorem{corollary}[theorem]{Corollary}
\newtheorem{remark}[theorem]{Remark}
\newtheorem{definition}[theorem]{Definition}
\newtheorem{example}[theorem]{Example}

\newcommand{\R}{\mathbb{R}}
\newcommand{\Q}{\mathbb{Q}}
\newcommand{\N}{\mathbb{N}}
\newcommand{\upf}{\operatorname{up}}
\newcommand{\sinc}{\operatorname{sinc}}
\newcommand{\eps}{\varepsilon}
\newcommand{\E}{\mathbb{E}}
\newcommand{\Prob}{\mathbb{P}}
\newcommand{\dd}{\,\mathrm{d}}
\newcommand{\pospart}[1]{\left(#1\right)_{+}}
\newcommand{\qbinom}[3]{\genfrac{[}{]}{0pt}{}{#1}{#2}_{#3}}
\newcommand{\poch}[3]{\left(#1;#2\right)_{#3}}
\newcommand{\G}{\mathcal{G}}
\newcommand{\F}{\mathcal{F}}
\newcommand{\depth}{\operatorname{depth}_2}

\title{\bfseries Exact Quarter-Base Extrapolation for Dyadic Samples\\
       of Rvachev's Up-Function Spline Prefixes}
\author{A technical note accompanying the ProveIt Fabius-function development}
\date{2 September 2026}

\begin{document}
\maketitle

\begin{abstract}
Let
\[
 \widehat{\upf}(t)=\prod_{k=0}^{\infty}
 \sinc\!\left(\frac{\pi t}{2^k}\right),
 \qquad \sinc z=\frac{\sin z}{z},
\]
and let $\upf_n$ be the inverse Fourier transform of the prefix through $k=n$.
For a fixed dyadic $x\in[-1,1]$, put $q_n=\upf_n(x)$.  We prove that after the
spline grid resolves $x$, the sequence is not merely asymptotic to a finite
sum of exponentials: it is \emph{exactly} of the form
\[
 q_n=\upf(x)+\sum_{r=1}^{d}A_r(x)4^{-rn},
 \qquad n\ge s,
\]
where $a=1-|x|=m/2^s$ is reduced and $d=\lfloor s/2\rfloor$.  For every
nontrivial depth $s\ge2$, all $A_1(x),\ldots,A_d(x)$ are nonzero, so this is
the exact number of geometric modes.

The constant term can consequently be recovered from only $d+1$ consecutive
spline values.  With $\rho=1/4$, any $N\ge s$ gives the single closed formula
\[
 \boxed{
 \upf(x)=\frac{1}{\poch{\rho}{\rho}{d}}
 \sum_{j=0}^{d}(-1)^j\rho^{\binom{j+1}{2}}
 \qbinom{d}{j}{\rho}\,q_{N+d-j}.}
\]
This is a quarter-base Gaussian-binomial/Richardson row.  Its generating
polynomial is $\poch{\rho z}{\rho}{d}/\poch{\rho}{\rho}{d}$, whose roots
annihilate all the decaying modes.  We also give the exact Vandermonde system
for all coefficients, a triangular extraction algorithm, stability bounds,
and exact rational examples.  The accompanying Python program uses only the
standard library and verifies every displayed numerical identity with
\texttt{fractions.Fraction} arithmetic.
\end{abstract}

\tableofcontents

\section{Statement of the problem and the answer}

We use the Fourier convention
\begin{equation}
 \widehat f(t)=\int_{\R}f(x)e^{-2\pi ixt}\dd x,
 \qquad
 f(x)=\int_{\R}\widehat f(t)e^{2\pi ixt}\dd t.
 \label{eq:fourier-convention}
\end{equation}
The Rvachev up-function is the even, nonnegative, compactly supported
$C^\infty$ function normalized by
\begin{equation}
 \widehat{\upf}(t)=\prod_{k=0}^{\infty}
 \sinc\!\left(\frac{\pi t}{2^k}\right).
 \label{eq:up-fourier-product}
\end{equation}
For $n\ge0$, define the finite-product approximation by
\begin{equation}
 \widehat{\upf_n}(t)=\prod_{k=0}^{n}
 \sinc\!\left(\frac{\pi t}{2^k}\right).
 \label{eq:finite-product-definition}
\end{equation}
There are $n+1$ sinc factors in \eqref{eq:finite-product-definition}; this
indexing point matters.  The inverse transform $\upf_n$ is a compactly
supported piecewise polynomial of degree $n$.

Fix a dyadic $x\in[-1,1]$ and put
\begin{equation}
 q_n:=\upf_n(x),
 \qquad
 a:=1-|x|\in[0,1].
 \label{eq:q-and-a}
\end{equation}
Write $a$ in reduced dyadic form
\begin{equation}
 a=\frac{m}{2^s},
 \qquad 0\le m\le2^s,
 \label{eq:reduced-a}
\end{equation}
where $s=0$ for $a\in\{0,1\}$ and, otherwise, $m$ is odd.  We call $s$ the
reduced dyadic depth of $a$ and set
\begin{equation}
 d:=\left\lfloor\frac{s}{2}\right\rfloor,
 \qquad \rho:=\frac14.
 \label{eq:d-and-rho}
\end{equation}

The main conclusions are the following.

\begin{theorem}[Exact eventual geometric tail]
\label{thm:main-tail}
For every dyadic $x\in[-1,1]$, with $s,d$ as above, there are rational
coefficients $A_1(x),\ldots,A_d(x)$ such that
\begin{equation}
 \boxed{
 q_n=\upf(x)+\sum_{k=1}^{d}A_k(x)\rho^{kn}
 \qquad(n\ge s).}
 \label{eq:main-tail}
\end{equation}
At the three depth-zero points, the statement is understood directly and is
valid already at $n=0$.  If $s\ge2$, then every displayed coefficient $A_k(x)$ is nonzero; hence the
number $d$ of decaying modes is exact.
\end{theorem}

\begin{theorem}[Single q-binomial extraction formula]
\label{thm:main-extraction}
Let $N\ge s$.  Then
\begin{equation}
 \boxed{
 \upf(x)=\frac{1}{\poch{\rho}{\rho}{d}}
 \sum_{j=0}^{d}(-1)^j\rho^{\binom{j+1}{2}}
 \qbinom{d}{j}{\rho}\,q_{N+d-j},
 \qquad \rho=\frac14.}
 \label{eq:main-extraction}
\end{equation}
Here
\begin{equation}
 \poch{z}{q}{d}=\prod_{r=0}^{d-1}(1-zq^r),
 \qquad
 \qbinom{d}{j}{q}
 =\frac{\poch{q}{q}{d}}
 {\poch{q}{q}{j}\poch{q}{q}{d-j}}.
 \label{eq:q-notation}
\end{equation}
In particular, the canonical choice $N=s$ recovers the exact rational value
from $q_s,q_{s+1},\ldots,q_{s+d}$.
\end{theorem}

The rest of the article proves these statements, identifies the coefficients,
and turns the closed formula into an exact algorithm.

\section{Finite sinc products are centered Fabius splines}

\subsection{Box densities and a random-sum interpretation}

For $k\ge0$, put
\begin{equation}
 b_k(x)=2^k\mathbf 1_{[-2^{-k-1},\,2^{-k-1}]}(x).
 \label{eq:box-density}
\end{equation}
A direct integration using \eqref{eq:fourier-convention} gives
\begin{equation}
 \widehat b_k(t)=\sinc\!\left(\frac{\pi t}{2^k}\right).
 \label{eq:box-transform}
\end{equation}
Consequently
\begin{equation}
 \upf_n=b_0*b_1*\cdots*b_n.
 \label{eq:finite-convolution}
\end{equation}
Equivalently, $\upf_n$ is the probability density of
\begin{equation}
 Z_n=\sum_{k=0}^{n}2^{-k}V_k,
 \qquad V_k\sim\operatorname{Unif}\!\left[-\frac12,\frac12\right]
 \quad\text{independently}.
 \label{eq:Zn}
\end{equation}
This immediately proves evenness and nonnegativity.  It also shows that the
support is
\begin{equation}
 \operatorname{supp}\upf_n
 =\left[-1+2^{-n-1},\,1-2^{-n-1}\right].
 \label{eq:finite-support}
\end{equation}

\subsection{The explicit Thue--Morse spline}

Let
\begin{equation}
 \eps_r=(-1)^{s_2(r)},
 \label{eq:thue-morse}
\end{equation}
where $s_2(r)$ is the number of $1$'s in the binary expansion of $r$.  For
$n\ge1$ define the centered spline
\begin{equation}
 S_n(y)=\frac{1}{2^{\binom n2}n!}
 \sum_{r\ge0}\eps_r\pospart{2^ny-\frac12-r}^{n}.
 \label{eq:centered-spline}
\end{equation}
At a fixed $y$, the sum is finite.  A convenient cutoff is
\begin{equation}
 N_n(y)=\max\!\left(0,\left\lfloor2^ny+\frac12\right\rfloor\right),
 \qquad 0\le r<N_n(y).
 \label{eq:spline-cutoff}
\end{equation}

\begin{lemma}[Exact indexing bridge]
\label{lem:indexing-bridge}
For every $n\ge1$ and every $x\in[-1,1]$,
\begin{equation}
 \boxed{\upf_n(x)=S_n(1+x).}
 \label{eq:upn-Sn}
\end{equation}
On the same interval, evenness also gives
\begin{equation}
 \upf_n(x)=S_n(1-|x|).
 \label{eq:upn-Sn-even}
\end{equation}
\end{lemma}

\begin{proof}
The standard truncated-power formula for the density of a sum of independent
uniform variables, applied to the widths $2^{-k}$ for $0\le k\le n$, gives
\begin{align}
 \upf_n(x)
 &=\frac{1}{n!\prod_{k=0}^{n}2^{-k}}
 \sum_{A\subseteq\{0,\ldots,n\}}(-1)^{|A|}
 \pospart{x+\frac12\sum_{k=0}^{n}2^{-k}-\sum_{k\in A}2^{-k}}^n.
 \label{eq:subset-box-spline}
\end{align}
Encode a subset $A$ by
\begin{equation}
 r=\sum_{k\in A}2^{n-k}.
\end{equation}
Then $r$ runs from $0$ to $2^{n+1}-1$ and $(-1)^{|A|}=\eps_r$.  Since
\begin{equation}
 \frac12\sum_{k=0}^{n}2^{-k}=1-2^{-n-1},
 \qquad
 \prod_{k=0}^{n}2^{-k}=2^{-n(n+1)/2},
\end{equation}
For $x\in[-1,1]$, no index beyond $2^{n+1}-1$ can contribute.  Hence
rescaling the positive part in \eqref{eq:subset-box-spline} yields
\begin{equation}
 \upf_n(x)=\frac{1}{2^{\binom n2}n!}
 \sum_{r\ge0}\eps_r
 \pospart{2^n(1+x)-\frac12-r}^{n},
\end{equation}
which is \eqref{eq:upn-Sn}.  Equation \eqref{eq:upn-Sn-even} follows because
every factor $b_k$ in \eqref{eq:finite-convolution} is even.
\end{proof}

\begin{remark}[The indexing checkpoint]
Some treatments use a prefix with factors $k=0,\ldots,n-1$.  That convention
produces $S_{n-1}$ rather than $S_n$.  In this note, the user's product ends at
$k=n$, contains $n+1$ factors, and produces the degree-$n$ spline $S_n$.
\end{remark}

\subsection{The Fabius distribution and the fold to up}

Let $U_1,U_2,\ldots$ be independent uniform random variables on $[0,1]$ and
put
\begin{equation}
 Y=\sum_{j=1}^{\infty}\frac{U_j}{2^j}.
 \label{eq:Y}
\end{equation}
Let $F$ be the distribution function of $Y$.  Then $F$ is the bounded Fabius
function \cite{fabius1966,arias1982,proveit}, and
\begin{equation}
 Z=2Y-1
 \label{eq:Z}
\end{equation}
has density $\upf$ on $[-1,1]$.  The familiar fold relation is
\begin{equation}
 \upf(x)=F(1-|x|),
 \qquad -1\le x\le1.
 \label{eq:fold}
\end{equation}
For
\begin{equation}
 Y_n=\sum_{j=1}^{n}\frac{U_j}{2^j},
 \qquad X_n=Y_n+2^{-n-1},
 \label{eq:Yn-Xn}
\end{equation}
the spline in \eqref{eq:centered-spline} is exactly the CDF of $X_n$:
\begin{equation}
 S_n(y)=\Prob(X_n\le y),
 \qquad 0\le y\le1.
 \label{eq:Sn-cdf}
\end{equation}
This follows either from the same uniform-sum truncated-power formula or by
induction from the smoothing recurrence.

Combining \eqref{eq:q-and-a}, \eqref{eq:upn-Sn-even}, and \eqref{eq:fold}, our
problem is therefore equivalent to this one:
\begin{equation}
 q_n=S_n(a)\longrightarrow F(a),
 \qquad a=\frac{m}{2^s}\in[0,1].
 \label{eq:reduced-problem}
\end{equation}

\section{Exact matched-cell moment deconvolution}

The decisive observation is that the omitted random tail has exactly the
half-width of one spline cell.

\subsection{The centered tail and its moments}

Let $Y'$ be an independent copy of $Y$.  The tail after $n$ binary scales is
\begin{equation}
 Y-Y_n\ \overset{\mathrm d}=\ 2^{-n}Y'.
\end{equation}
Consequently, with $Z'=2Y'-1$ and
\begin{equation}
 \delta_n:=2^{-n-1},
 \label{eq:delta}
\end{equation}
we have the centered decomposition
\begin{equation}
 Y\ \overset{\mathrm d}=\ X_n+\delta_n Z'.
 \label{eq:centered-tail-decomposition}
\end{equation}
Conditioning on $Z'$ gives the exact convolution identity
\begin{equation}
 F(t)=\E\bigl[S_n(t-\delta_n Z)\bigr]
 =\int_{-1}^{1}S_n(t-\delta_n z)\upf(z)\dd z.
 \label{eq:F-Sn-convolution}
\end{equation}

Define the even moments
\begin{equation}
 c_k:=\E(Z^{2k})=\int_{-1}^{1}z^{2k}\upf(z)\dd z,
 \qquad c_0=1.
 \label{eq:ck}
\end{equation}
The centered moment-generating function is
\begin{align}
 \mathscr C(T)
 &:=\E(e^{TZ})
   =\sum_{k\ge0}c_k\frac{T^{2k}}{(2k)!} \\
 &=\prod_{j=1}^{\infty}
   \frac{\sinh(T/2^j)}{T/2^j}.
 \label{eq:C-product}
\end{align}
The product is the real-exponential counterpart of the sinc product
\eqref{eq:up-fourier-product}.

Define reciprocal centered moments $\gamma_k$ by
\begin{equation}
 \frac{1}{\mathscr C(T)}
 =\sum_{k\ge0}\gamma_k\frac{T^{2k}}{(2k)!}.
 \label{eq:gamma-definition}
\end{equation}
Thus
\begin{equation}
 \gamma_0=1,
 \qquad
 \sum_{j=0}^{r}\binom{2r}{2j}\gamma_jc_{r-j}=0
 \quad(r\ge1).
 \label{eq:gamma-convolution}
\end{equation}
The first values are
\begin{equation}
 \gamma_0=1,
 \quad \gamma_1=-\frac19,
 \quad \gamma_2=\frac{31}{675},
 \quad \gamma_3=-\frac{2347}{59535},
 \quad \gamma_4=\frac{1904369}{32531625}.
 \label{eq:gamma-first}
\end{equation}

\begin{lemma}[Rationality and nonvanishing of the reciprocal moments]
\label{lem:gamma-sign}
Every $\gamma_k$ is rational, and
\begin{equation}
 (-1)^k\gamma_k>0.
 \label{eq:gamma-sign}
\end{equation}
In particular, $\gamma_k\ne0$ for every $k$.
\end{lemma}

\begin{proof}
The logarithm of \eqref{eq:C-product} is
\begin{equation}
 \log\mathscr C(T)
 =\sum_{r\ge1}
 \frac{2^{2r-1}B_{2r}}
 {r(2r)!(2^{2r}-1)}T^{2r},
 \label{eq:log-C}
\end{equation}
so its coefficients, and hence those of its exponential reciprocal, are
rational.  For the sign, use
\begin{equation}
 \frac{z}{\sinh z}
 =\sum_{r\ge0}(-1)^r b_r z^{2r},
 \qquad b_r>0.
 \label{eq:z-over-sinh-sign}
\end{equation}
For $r\ge1$, this follows explicitly from
\begin{equation}
 [z^{2r}]\frac{z}{\sinh z}
 =\frac{2(1-2^{2r-1})B_{2r}}{(2r)!}.
\end{equation}
Now
\begin{equation}
 \frac1{\mathscr C(T)}
 =\prod_{j=1}^{\infty}\frac{T/2^j}{\sinh(T/2^j)}.
\end{equation}
Every contribution to the coefficient of $T^{2k}$ has the common sign
$(-1)^k$, and at least one contribution is strictly positive in absolute
value.  This proves \eqref{eq:gamma-sign}.
\end{proof}

\subsection{A dyadic grid point is the center of one polynomial cell}

The knots of $S_n$ are
\begin{equation}
 \kappa_{r,n}=\frac{r+1/2}{2^n}.
 \label{eq:knots}
\end{equation}
If $a=M/2^n$, then
\begin{equation}
 [a-\delta_n,a+\delta_n]
 =\left[\frac{M-1/2}{2^n},\frac{M+1/2}{2^n}\right]
 \label{eq:matched-cell}
\end{equation}
is exactly one closed spline cell, centered at $a$.  On its interior, $S_n$
is a single polynomial of degree at most $n$.  The exterior cells at $a=0$
and $a=1$ are constant and cause no exception.

\begin{theorem}[Exact reciprocal moment deconvolution on a matched cell]
\label{thm:matched-deconvolution}
Let $n\ge1$ and let $a=M/2^n\in[0,1]$.  Then
\begin{equation}
 \boxed{
 S_n(a)=\sum_{k=0}^{\lfloor n/2\rfloor}
 \frac{\gamma_k}{(2k)!}\,4^{-k(n+1)}F^{(2k)}(a).}
 \label{eq:matched-deconvolution}
\end{equation}
More generally, for $0\le r\le n$, the $n$-jet satisfies the triangular
moment relation
\begin{equation}
 F^{(r)}(a)=
 \sum_{k=0}^{\lfloor(n-r)/2\rfloor}
 \frac{c_k\delta_n^{2k}}{(2k)!}
 S_n^{(r+2k)}(a).
 \label{eq:jet-moment-relation}
\end{equation}
\end{theorem}

\begin{proof}
Let $P(h)=S_n(a+h)$ denote the polynomial branch on
$[-\delta_n,\delta_n]$.  Because $a-\delta_n Z$ remains in this cell almost
surely, differentiating the compactly supported convolution
\eqref{eq:F-Sn-convolution} in the distributional sense, and then using the
ordinary polynomial branch inside the cell, gives
\begin{equation}
 F^{(r)}(a)=\E\bigl[P^{(r)}(-\delta_n Z)\bigr]
 \qquad(0\le r\le n).
 \label{eq:derivative-expectation}
\end{equation}
The Taylor formula for the polynomial $P^{(r)}$ is finite.  Odd moments of $Z$
vanish because $\upf$ is even, so taking expectations gives
\eqref{eq:jet-moment-relation}.

To invert the triangular relation, consider
\begin{equation}
 T:=\sum_{k=0}^{\lfloor n/2\rfloor}
 \frac{\gamma_k\delta_n^{2k}}{(2k)!}F^{(2k)}(a).
\end{equation}
Substitute \eqref{eq:jet-moment-relation}.  The coefficient of
$\delta_n^{2r}S_n^{(2r)}(a)/(2r)!$ becomes
\begin{equation}
 \sum_{k=0}^{r}\binom{2r}{2k}\gamma_kc_{r-k}.
\end{equation}
By \eqref{eq:gamma-convolution}, this coefficient is $1$ for $r=0$ and $0$
for every $r\ge1$.  Therefore $T=S_n(a)$.  Finally,
$\delta_n^{2k}=2^{-2k(n+1)}=4^{-k(n+1)}$, proving
\eqref{eq:matched-deconvolution}.
\end{proof}

\begin{remark}[Why the formula is exact]
There is no passage to an infinite Taylor series in this proof.  At a matched
dyadic point, the tail displacement $\delta_nZ$ stays inside one degree-$n$
polynomial cell.  Taylor expansion therefore terminates, and reciprocal
moment deconvolution is finite algebra.  This is the mechanism behind the
eventual exact geometric behavior.
\end{remark}

\section{Dyadic depth kills all but finitely many derivative modes}

\subsection{The signed global extension and derivative dilation}

Introduce the signed global extension
\begin{equation}
 \G(y)=\sum_{b\ge0}\eps_b\upf(y-2b-1).
 \label{eq:global-extension}
\end{equation}
The sum is locally finite.  On the first block,
\begin{equation}
 \G(y)=\upf(y-1),\qquad 0\le y\le2,
 \label{eq:first-block}
\end{equation}
and on $[0,1]$ it agrees with $F$.  It satisfies the global differential
refinement equation \cite{arias2018,proveit}
\begin{equation}
 \G'(y)=2\G(2y).
 \label{eq:G-diffeq}
\end{equation}
Iterating gives
\begin{equation}
 \G^{(r)}(y)=2^{\binom{r+1}{2}}\G(2^ry).
 \label{eq:G-derivatives}
\end{equation}
In particular, for $a\in[0,1]$,
\begin{equation}
 F^{(2k)}(a)=2^{k(2k+1)}\G(4^ka).
 \label{eq:F-even-derivative}
\end{equation}
The block endpoints and centers give
\begin{align}
 \G(2b)&=0, \\
 \G(2b+1)&=\eps_b, \\
 \G\!\left(2b+\frac12\right)
 =\G\!\left(2b+\frac32\right)&=\frac{\eps_b}{2}.
 \label{eq:G-integer-half-values}
\end{align}

\subsection{The finite derivative jet at a dyadic point}

Let $a=m/2^s\in(0,1)$ be reduced, so $m$ is odd, and put
$d=\lfloor s/2\rfloor$.

\begin{lemma}[Vanishing beyond half the dyadic depth]
\label{lem:derivative-vanishing}
For every $k>d$,
\begin{equation}
 F^{(2k)}(a)=0.
 \label{eq:derivative-vanishing}
\end{equation}
For every $0\le k\le d$, one has
\begin{equation}
 F^{(2k)}(a)\ne0.
 \label{eq:surviving-derivatives-nonzero}
\end{equation}
\end{lemma}

\begin{proof}
If $s=2d$, then $4^da=m$ is an odd integer, while for every $k>d$,
$4^ka=m4^{k-d}$ is an even integer.  If $s=2d+1$, then $4^da=m/2$ is an odd
half-integer, while for $k>d$, $4^ka=2m4^{k-d-1}$ is an even integer.  Now use
\eqref{eq:F-even-derivative} and \eqref{eq:G-integer-half-values}.  At the top
index, $\G(4^da)$ is respectively $\pm1$ or $\pm1/2$.  For $k<d$, the
number $4^ka=m/2^{s-2k}$ is not an integer, hence it lies in the interior of
one of the two-unit blocks in \eqref{eq:global-extension}; the up-function is
strictly positive in the interior of its support, so $\G(4^ka)\ne0$.
\end{proof}

\subsection{Proof of the exact eventual tail}

For $n\ge s$, the reduced dyadic $a=m/2^s$ is also a point of the $2^{-n}$
grid:
\begin{equation}
 a=\frac{m2^{n-s}}{2^n}.
\end{equation}
Apply Theorem~\ref{thm:matched-deconvolution}, then use
Lemma~\ref{lem:derivative-vanishing}.  We obtain
\begin{align}
 q_n=S_n(a)
 &=F(a)+\sum_{k=1}^{d}
 \frac{\gamma_k}{(2k)!}4^{-k(n+1)}F^{(2k)}(a) \\
 &=\upf(x)+\sum_{k=1}^{d}A_k(x)4^{-kn},
 \label{eq:tail-derived}
\end{align}
where
\begin{equation}
 \boxed{
 A_k(x)=\frac{\gamma_k}{(2k)!4^k}F^{(2k)}(a)
 =\frac{\gamma_k2^{k(2k-1)}}{(2k)!}\G(4^ka),
 \qquad a=1-|x|.}
 \label{eq:Ak}
\end{equation}
This proves the formula in Theorem~\ref{thm:main-tail}.  It also proves
rationality without requiring a separate arithmetic theorem about $F$ or
$\G$.  Indeed, each $q_n=S_n(a)$ is rational by the finite sum
\eqref{eq:centered-spline}.  Evaluating \eqref{eq:tail-derived} at
$n=s,s+1,\ldots,s+d$ gives a square Vandermonde system over $\Q$ for
$F(a),A_1(x),\ldots,A_d(x)$, with distinct nodes
$1,\rho,\ldots,\rho^d$.  Its unique solution is rational.  In particular,
this argument itself recovers the classical rationality of $\upf(x)=F(a)$ at
dyadic arguments.

For $s\ge2$, Lemma~\ref{lem:gamma-sign} and
Lemma~\ref{lem:derivative-vanishing} imply $A_k(x)\ne0$ for every
$1\le k\le d$.  More explicitly, if
$a=m/2^s$ is reduced, then
\begin{equation}
 A_d(x)=
 \begin{cases}
 \displaystyle
 \frac{\gamma_d2^{d(2d-1)}}{(2d)!}\,
 \eps_{(m-1)/2},&s=2d,\\[1.2em]
 \displaystyle
 \frac{\gamma_d2^{d(2d-1)-1}}{(2d)!}\,
 \eps_{\lfloor m/4\rfloor},&s=2d+1.
 \end{cases}
 \label{eq:top-coefficient}
\end{equation}
Thus the tail contains exactly $d$ nonconstant geometric modes, all with
nonzero coefficients.

\begin{remark}[What happens before $n=s$]
For $n<s$, the point $a$ is not a center of the $2^{-n}$ spline grid.  The
tail interval need not coincide with one polynomial cell, so the matched-cell
deconvolution argument does not apply in this form.  These are precisely the
``few initial values'' allowed in the question.  Once $n\ge s$, the formula
\eqref{eq:tail-derived} is an identity for every subsequent $n$.
\end{remark}

\section{Extraction from consecutive spline values}

We now solve the purely algebraic problem created by
Theorem~\ref{thm:main-tail}.  Let
\begin{equation}
 q_n=L+\sum_{r=1}^{d}A_r\rho^{rn},
 \qquad \rho=\frac14,
 \qquad n\ge N.
 \label{eq:abstract-tail}
\end{equation}
In the application, $L=\upf(x)$ and any $N\ge s$ is valid.

\subsection{The Vandermonde system}

The $d+1$ values $q_N,\ldots,q_{N+d}$ satisfy
\begin{equation}
 \begin{bmatrix}
 1&\rho^N&\rho^{2N}&\cdots&\rho^{dN}\\
 1&\rho^{N+1}&\rho^{2(N+1)}&\cdots&\rho^{d(N+1)}\\
 \vdots&\vdots&\vdots&&\vdots\\
 1&\rho^{N+d}&\rho^{2(N+d)}&\cdots&\rho^{d(N+d)}
 \end{bmatrix}
 \begin{bmatrix}L\\A_1\\A_2\\\vdots\\A_d\end{bmatrix}
 =
 \begin{bmatrix}q_N\\q_{N+1}\\\vdots\\q_{N+d}\end{bmatrix}.
 \label{eq:vandermonde-system}
\end{equation}
The determinant is nonzero because $1,\rho,\ldots,\rho^d$ are distinct.
Therefore all coefficients are uniquely determined over $\Q$.

A useful interpolation reformulation is
\begin{equation}
 P_N(z):=L+\sum_{r=1}^{d}A_r\rho^{rN}z^r.
 \label{eq:PN}
\end{equation}
Then
\begin{equation}
 P_N(\rho^j)=q_{N+j},
 \qquad 0\le j\le d,
 \label{eq:PN-samples}
\end{equation}
and the desired limit is simply
\begin{equation}
 L=P_N(0).
 \label{eq:L-P0}
\end{equation}
Moreover, once $P_N$ has been interpolated,
\begin{equation}
 A_r=\rho^{-rN}[z^r]P_N(z).
 \label{eq:all-Ar}
\end{equation}
Thus the same $d+1$ values recover every coefficient, not only the limit.
Explicitly,
\begin{equation}
 P_N(z)=\sum_{j=0}^{d}q_{N+j}
 \prod_{\substack{0\le\ell\le d\\\ell\ne j}}
 \frac{z-\rho^\ell}{\rho^j-\rho^\ell}.
 \label{eq:lagrange-P}
\end{equation}

\subsection{Lagrange weights at zero}

Setting $z=0$ in \eqref{eq:lagrange-P}, the weight of $q_{N+j}$ is
\begin{align}
 \lambda_{d,j}(\rho)
 &:=\prod_{\substack{0\le\ell\le d\\\ell\ne j}}
 \frac{-\rho^\ell}{\rho^j-\rho^\ell} \\
 &=(-1)^{d-j}
 \frac{\rho^{\binom{d-j+1}{2}}}
 {\poch{\rho}{\rho}{j}\poch{\rho}{\rho}{d-j}}.
 \label{eq:forward-lagrange-weight}
\end{align}
Hence a forward-order formula is
\begin{equation}
 L=\sum_{j=0}^{d}(-1)^{d-j}
 \frac{\rho^{\binom{d-j+1}{2}}}
 {\poch{\rho}{\rho}{j}\poch{\rho}{\rho}{d-j}}q_{N+j}.
 \label{eq:forward-extraction}
\end{equation}
Reindexing by $j\mapsto d-j$ and using the Gaussian binomial coefficient in
\eqref{eq:q-notation} gives exactly
\begin{equation}
 L=\frac{1}{\poch{\rho}{\rho}{d}}
 \sum_{j=0}^{d}(-1)^j\rho^{\binom{j+1}{2}}
 \qbinom{d}{j}{\rho}\,q_{N+d-j},
 \label{eq:qbinomial-extraction-again}
\end{equation}
which proves Theorem~\ref{thm:main-extraction}.

\subsection{The row polynomial and mode annihilation}

Define the reverse-order weights
\begin{equation}
 w_{d,j}=\frac{(-1)^j\rho^{\binom{j+1}{2}}}
 {\poch{\rho}{\rho}{d}}\qbinom{d}{j}{\rho},
 \qquad 0\le j\le d.
 \label{eq:wdj}
\end{equation}
The finite $q$-binomial theorem \cite{gasperrahman} gives their generating polynomial:
\begin{align}
 W_d(z)
 &:=\sum_{j=0}^{d}w_{d,j}z^j \\
 &=\frac{\poch{\rho z}{\rho}{d}}
 {\poch{\rho}{\rho}{d}}
 =\prod_{r=1}^{d}\frac{1-\rho^rz}{1-\rho^r}.
 \label{eq:row-polynomial}
\end{align}
Therefore
\begin{equation}
 W_d(1)=1,
 \qquad
 W_d(\rho^{-r})=0\quad(1\le r\le d).
 \label{eq:row-roots}
\end{equation}
In the reverse combination $\sum_jw_{d,j}q_{N+d-j}$, the constant mode
contributes $LW_d(1)=L$, while the mode $A_r\rho^{rn}$ contributes a multiple
of $W_d(\rho^{-r})=0$.  This is the shortest proof of the extraction formula
once the geometric tail is known.

\subsection{Shift-operator and triangular Richardson forms}

Let $E$ be the forward shift, $(Eq)_n=q_{n+1}$.  Since $E$ acts on the mode
$\rho^{rn}$ by the eigenvalue $\rho^r$, we have
\begin{equation}
 \boxed{
 L=\frac{\prod_{r=1}^{d}(E-\rho^r)}
 {\prod_{r=1}^{d}(1-\rho^r)}q_N.}
 \label{eq:shift-operator}
\end{equation}
Expanding \eqref{eq:shift-operator} is another route to
\eqref{eq:qbinomial-extraction-again}.

It also gives an in-place triangular algorithm.  Start with
\begin{equation}
 R_0(j)=q_{N+j},\qquad 0\le j\le d,
\end{equation}
and, for $r=1,\ldots,d$, form
\begin{equation}
 R_r(j)=\frac{R_{r-1}(j+1)-\rho^rR_{r-1}(j)}{1-\rho^r},
 \qquad 0\le j\le d-r.
 \label{eq:richardson-recursion}
\end{equation}
Then
\begin{equation}
 R_d(0)=L.
 \label{eq:richardson-result}
\end{equation}
With rational input, every operation is rational.

\subsection{Sample count and stability}

For $s\ge2$, the top coefficient $A_d$ is nonzero.  Within the model
\eqref{eq:abstract-tail}, there are $d+1$ independent unknowns
$L,A_1,\ldots,A_d$, and the geometric Vandermonde matrix is nonsingular.
Thus $d+1$ samples are both sufficient and, without additional information,
necessary to determine all coefficients.

For approximate rather than exact inputs, the extraction row is well
conditioned.  Since $0<\rho<1$ and the signs in \eqref{eq:wdj} alternate,
\begin{equation}
 \sum_{j=0}^{d}|w_{d,j}|
 =W_d(-1)
 =\frac{\poch{-\rho}{\rho}{d}}{\poch{\rho}{\rho}{d}}.
 \label{eq:row-variation}
\end{equation}
At $\rho=1/4$, this increases to approximately
\begin{equation}
 \frac{\poch{-1/4}{1/4}{\infty}}
 {\poch{1/4}{1/4}{\infty}}
 \approx1.969260354.
 \label{eq:variation-numeric}
\end{equation}
Hence if every supplied $q_n$ has absolute error at most $\eta$, the extracted
constant has error at most about $1.97\eta$, uniformly in the depth.  Exact
rational arithmetic, as used in the accompanying program, eliminates this
issue entirely.

\subsection{First extraction rows}

For convenience, the first quarter-base rows are
\begin{align}
 d=0:\quad &L=q_N,\\
 d=1:\quad &L=\frac43q_{N+1}-\frac13q_N,\\
 d=2:\quad &L=\frac{64}{45}q_{N+2}-\frac49q_{N+1}+\frac1{45}q_N,\\
 d=3:\quad &L=\frac{4096}{2835}q_{N+3}-\frac{64}{135}q_{N+2}
             +\frac4{135}q_{N+1}-\frac1{2835}q_N.
 \label{eq:first-rows}
\end{align}

\section{Worked exact examples}

\subsection{Depth two: \texorpdfstring{$x=\pm3/4$}{x=+/-3/4}}

Here
\begin{equation}
 a=1-|x|=\frac14,
 \qquad s=2,
 \qquad d=1.
\end{equation}
The exact dyadic value and derivative are
\begin{equation}
 F\!\left(\frac14\right)=\frac5{72},
 \qquad
 F''\!\left(\frac14\right)=8\G(1)=8.
\end{equation}
Since $\gamma_1=-1/9$, formula \eqref{eq:Ak} gives $A_1=-1/9$.  Therefore
\begin{equation}
 \boxed{
 \upf_n\!\left(\pm\frac34\right)
 =\frac5{72}-\frac19\,4^{-n},
 \qquad n\ge2.}
 \label{eq:quarter-example-tail}
\end{equation}
The first two values in the valid range are
\begin{equation}
 q_2=\frac1{16},
 \qquad
 q_3=\frac{13}{192}.
\end{equation}
The $d=1$ extraction row gives
\begin{equation}
 \frac43q_3-\frac13q_2
 =\frac{4q_3-q_2}{3}
 =\frac5{72}=\upf\!\left(\pm\frac34\right).
\end{equation}

\subsection{Depth four: \texorpdfstring{$x=\pm11/16$}{x=+/-11/16}}

Now
\begin{equation}
 a=1-|x|=\frac5{16},
 \qquad s=4,
 \qquad d=2.
\end{equation}
The exact value is
\begin{equation}
 L=F\!\left(\frac5{16}\right)
 =\frac{305857}{2073600}.
 \label{eq:five-sixteenths-value}
\end{equation}
The derivative dilation gives
\begin{align}
 F''\!\left(\frac5{16}\right)
 &=8\G\!\left(\frac54\right)
 =8\upf\!\left(\frac14\right)
 =8F\!\left(\frac34\right)
 =\frac{67}{9},\\
 F^{(4)}\!\left(\frac5{16}\right)
 &=2^{10}\G(5)=-1024.
 \label{eq:five-sixteenths-derivatives}
\end{align}
Using $\gamma_1=-1/9$ and $\gamma_2=31/675$ in \eqref{eq:Ak},
\begin{equation}
 A_1=-\frac{67}{648},
 \qquad
 A_2=-\frac{248}{2025}.
\end{equation}
Thus the eventual formula is
\begin{equation}
 \boxed{
 \upf_n\!\left(\pm\frac{11}{16}\right)
 =\frac{305857}{2073600}
 -\frac{67}{648}\,4^{-n}
 -\frac{248}{2025}\,16^{-n},
 \qquad n\ge4.}
 \label{eq:five-sixteenths-tail}
\end{equation}
The exact spline values are
\begin{align}
 q_4&=\frac{1205}{8192},\\
 q_5&=\frac{289799}{1966080},\\
 q_6&=\frac{13917509}{94371840}.
 \label{eq:q456}
\end{align}
Applying the $d=2$ row,
\begin{equation}
 \frac{64}{45}q_6-\frac49q_5+\frac1{45}q_4
 =\frac{305857}{2073600},
 \label{eq:example-extraction}
\end{equation}
exactly as claimed.

\subsection{Depth zero and one}

When $d=0$, there is no decaying term.  The spline sequence stabilizes as soon
as the point lies on the grid.  In particular,
\begin{align}
 \upf_n(0)&=1,\\
 \upf_n(\pm1)&=0,\\
 \upf_n\!\left(\pm\frac12\right)&=\frac12\qquad(n\ge1).
\end{align}
These are the $s=0$ and $s=1$ cases of the main theorem.

\section{Exact computational method}

For a given dyadic $x\in[-1,1]$, the shortest exact workflow is:

\begin{enumerate}
 \item Reduce $a=1-|x|$ to $m/2^s$ and set $d=\lfloor s/2\rfloor$.
 \item Evaluate $q_s,q_{s+1},\ldots,q_{s+d}$ from the finite spline formula
       \eqref{eq:centered-spline}.  Every operation is rational.
 \item Apply either the single q-binomial formula
       \eqref{eq:main-extraction} or the triangular recursion
       \eqref{eq:richardson-recursion}.  The result is exactly $\upf(x)$.
 \item To recover the complete eventual formula, solve the rational
       Vandermonde system \eqref{eq:vandermonde-system}, or interpolate
       $P_s$ using \eqref{eq:lagrange-P} and use \eqref{eq:all-Ar}.
 \item Optionally verify the recovered coefficients against one or more later
       values.  The theorem guarantees exact agreement for every $n\ge s$.
\end{enumerate}

The included file \texttt{verify\_extrapolation.py} implements all of these
steps with Python's \texttt{Fraction} type.  It contains:

\begin{itemize}
 \item exact Thue--Morse spline evaluation;
 \item dyadic-depth detection;
 \item finite q-Pochhammer and Gaussian-binomial functions;
 \item both closed-row and recursive extraction;
 \item exact Gauss--Jordan solution of the full Vandermonde system;
 \item computation of reciprocal moments $\gamma_k$ from Bernoulli numbers;
 \item assertions for the two examples above and checks at later $n$;
 \item an optional exhaustive exact test over all reduced dyadics through a
       user-selected denominator depth.
\end{itemize}

A typical invocation is
\begin{lstlisting}[language=bash]
python3 verify_extrapolation.py
python3 verify_extrapolation.py --x=-11/16
python3 verify_extrapolation.py --self-test-max-depth 7
\end{lstlisting}
The direct Thue--Morse sum has a cutoff of order $2^n a$, so it is intended as
an exact verification tool at moderate depth.  Faster evaluation can use the
binary recurrences and formal arithmetic developed in the ProveIt sources;
the extrapolation row itself is only quadratic-time in $d$ in its recursive
form.

\section{Relation to the half-base q-row in the ProveIt development}

The linked development contains a general half-base Gaussian-binomial row for
shifted finite splines.  Its row polynomial has the form
\begin{equation}
 \frac{(z/2;1/2)_n}{(1/2;1/2)_n}
\end{equation}
and annihilates modes $2^{-rn}$.  The result here is tailored to the specific
centered prefix
\begin{equation}
 \prod_{k=0}^{n}\sinc\!\left(\frac{\pi t}{2^k}\right).
\end{equation}
The centered tail law is even, so every odd moment vanishes.  In the exact
matched-cell deconvolution, only even derivative orders $2k$ survive, and the
scale factor is
\begin{equation}
 2^{-2k(n+1)}=4^{-k(n+1)}.
\end{equation}
After absorbing the fixed factor $4^{-k}$ into $A_k$, the sequence modes are
$4^{-kn}$.  That is why the natural and minimal base in
\eqref{eq:main-extraction} is $\rho=1/4$, not $1/2$.  The quarter-base row uses
only $\lfloor s/2\rfloor+1$ consecutive values, whereas a generic half-base
filter that also prepares for absent odd modes would be longer.

There is also a notational distinction worth preserving: in parts of the
linked source, a letter $Q$ denotes a complex \emph{translation parameter} in
a shifted spline.  The number $\rho=1/4$ in this note is instead the
\emph{geometric decay ratio} of the centered finite-product sequence.

\section{Conclusion}

For a fixed dyadic input, the finite sinc-product approximations have a much
stronger structure than ordinary convergence.  Once the spline grid reaches
the reduced denominator depth, the error is an exact finite combination of
known geometric ratios:
\begin{equation}
 \upf_n(x)-\upf(x)
 =A_1(x)4^{-n}+A_2(x)4^{-2n}+\cdots
 +A_{\lfloor s/2\rfloor}(x)4^{-\lfloor s/2\rfloor n}.
\end{equation}
The analytic reason is the exact centering of the omitted random tail; the
arithmetic reason is the vanishing of high even derivatives at dyadic points;
and the algebraic reason the limit is extractable is the q-Vandermonde
structure of the remaining modes.

The normalized q-Pochhammer polynomial
\begin{equation}
 \frac{\poch{\rho z}{\rho}{d}}{\poch{\rho}{\rho}{d}}
\end{equation}
packages the entire extraction row.  Its value at $1$ preserves the true
constant, and its roots $\rho^{-1},\ldots,\rho^{-d}$ annihilate every decaying
mode.  Formula \eqref{eq:main-extraction} is therefore both a closed proof and
a practical exact algorithm for recovering the rational value $\upf(x)$ from
finitely many rational spline approximations.

\appendix

\section{A compact theorem-to-algorithm summary}

Let $x\in[-1,1]$ be dyadic.

\begin{align*}
 a&=1-|x|=m/2^s\quad\text{in lowest terms},\\
 d&=\lfloor s/2\rfloor,\qquad \rho=1/4,\\
 q_n&=\upf_n(x)=S_n(a).
\end{align*}
Then, for every $n\ge s$,
\begin{align*}
 q_n
 &=\upf(x)+\sum_{k=1}^{d}A_k(x)\rho^{kn},\\
 A_k(x)
 &=\frac{\gamma_k}{(2k)!4^k}F^{(2k)}(a),\\
 \frac1{\mathscr C(T)}
 &=\sum_{k\ge0}\gamma_k\frac{T^{2k}}{(2k)!},\\
 \mathscr C(T)
 &=\prod_{j=1}^{\infty}\frac{\sinh(T/2^j)}{T/2^j}.
\end{align*}
Finally, for any $N\ge s$,
\begin{equation*}
 \upf(x)=\frac{1}{\poch{1/4}{1/4}{d}}
 \sum_{j=0}^{d}(-1)^j4^{-\binom{j+1}{2}}
 \qbinom{d}{j}{1/4}\upf_{N+d-j}(x).
\end{equation*}
This one line is the requested extraction formula.

\begin{thebibliography}{99}

\bibitem{fabius1966}
J.~Fabius,
\newblock A probabilistic example of a nowhere analytic $C^\infty$-function,
\newblock \emph{Zeitschrift f\"ur Wahrscheinlichkeitstheorie und Verwandte
Gebiete} \textbf{5} (1966), 173--174.

\bibitem{rvachev1990}
V.~A.~Rvachev,
\newblock Compactly supported solutions of functional-differential equations
and their applications,
\newblock \emph{Russian Mathematical Surveys} \textbf{45} (1990), no.~1,
87--120.

\bibitem{arias1982}
J.~Arias de Reyna,
\newblock Definici\'on y estudio de una funci\'on indefinidamente diferenciable
de soporte compacto,
\newblock \emph{Revista de la Real Academia de Ciencias Exactas, F\'isicas y
Naturales de Madrid} \textbf{76} (1982), no.~1, 21--38;
English translation, arXiv:1702.05442 (2017).

\bibitem{arias2018}
J.~Arias de Reyna,
\newblock Arithmetic of the Fabius function,
\newblock \emph{INTEGERS} \textbf{18} (2018), Article A51;
arXiv:1702.06487, version~3.

\bibitem{gasperrahman}
G.~Gasper and M.~Rahman,
\newblock \emph{Basic Hypergeometric Series}, second edition,
\newblock Cambridge University Press, 2004.

\bibitem{proveit}
V.~Reshetnikov,
\newblock \emph{ProveIt: Analysis/FabiusFunction}, Lean~4 formal development
and accompanying articles,
\newblock \url{https://github.com/VladimirReshetnikov/ProveIt/tree/main/Analysis/FabiusFunction},
consulted 2 September 2026.

\end{thebibliography}

\end{document}
